\documentclass{beamer}
\usepackage[utf8]{inputenc}

\usetheme{Madrid}
\usecolortheme{default}
\usepackage{amsmath,amssymb,amsfonts,amsthm}
\usepackage{txfonts}
\usepackage{tkz-euclide}
\usepackage{listings}
\usepackage{adjustbox}
\usepackage{array}
\usepackage{tabularx}
\usepackage{gvv}
\usepackage{lmodern}
\usepackage{circuitikz}
\usepackage{tikz}
\usepackage{graphicx}

\setbeamertemplate{page number in head/foot}[totalframenumber]

\usepackage{tcolorbox}
\tcbuselibrary{minted,breakable,xparse,skins}



\definecolor{bg}{gray}{0.95}
\DeclareTCBListing{mintedbox}{O{}m!O{}}{%
  breakable=true,
  listing engine=minted,
  listing only,
  minted language=#2,
  minted style=default,
  minted options={%
    linenos,
    gobble=0,
    breaklines=true,
    breakafter=,,
    fontsize=\small,
    numbersep=8pt,
    #1},
  boxsep=0pt,
  left skip=0pt,
  right skip=0pt,
  left=25pt,
  right=0pt,
  top=3pt,
  bottom=3pt,
  arc=5pt,
  leftrule=0pt,
  rightrule=0pt,
  bottomrule=2pt,
  toprule=2pt,
  colback=bg,
  colframe=orange!70,
  enhanced,
  overlay={%
    \begin{tcbclipinterior}
    \fill[orange!20!white] (frame.south west) rectangle ([xshift=20pt]frame.north west);
    \end{tcbclipinterior}},
  #3,
}
\lstset{
    language=C,
    basicstyle=\ttfamily\small,
    keywordstyle=\color{blue},
    stringstyle=\color{orange},
    commentstyle=\color{green!60!black},
    numbers=left,
    numberstyle=\tiny\color{gray},
    breaklines=true,
    showstringspaces=false,
}
%------------------------------------------------------------

\title
{12.427}
\author 
{AI25BTECH11034 - Sujal Chauhan }



\begin{document}

\frame{\titlepage}
\begin{frame}{Question}
    Let $\mathbf{M}$ be a $3 \times 3$ matrix such that
\[
\mathbf{M} \begin{pmatrix} -2 \\ 1 \\ 0 \end{pmatrix} = \begin{pmatrix} 6 \\ -3 \\ 0 \end{pmatrix}
\]
and suppose that
\[
\mathbf{M}^3 \begin{pmatrix} 1 \\ -\frac{1}{2} \\ 0 \end{pmatrix} = \begin{pmatrix} \alpha \\ \beta \\ \gamma \end{pmatrix}
\]
for some $\alpha, \beta, \gamma \in \mathbb{R}$. Then $|\alpha|$ is equal to \rule{2cm}{0.15mm}.

\end{frame}

\begin{frame}{Solution}
    Let $\Vec{M}$ be a $3\times3$ matrix such that
    \begin{align}
        \Vec{M}\myvec{-2\\1\\0} = \myvec{6\\-3\\0}
    \end{align}
    We observe that
    \begin{align}
        \myvec{6\\-3\\0} = -3\myvec{-2\\1\\0}
    \end{align}
    Hence, $\myvec{-2\\1\\0}$ is an eigenvector of $\Vec{M}$ with eigenvalue $\lambda = -3$.

    \vspace{0.4cm}
    \end{frame}
    \begin{frame}{Solution}
    Now consider
    \begin{align}
        \Vec{M}^3\myvec{1\\-\tfrac{1}{2}\\0} = \myvec{\alpha\\\beta\\\gamma}
    \end{align}
    Note that
    \begin{align}
        \myvec{1\\-\tfrac{1}{2}\\0} = -\tfrac{1}{2}\myvec{-2\\1\\0}
    \end{align}
    which is a scalar multiple of the same eigenvector.  
    Therefore, it also corresponds to the same eigenvalue $\lambda = -3$.

    \vspace{0.4cm}
    \end{frame}
    \begin{frame}{Solution}
        
    
    Applying the matrix $\Vec{M}$ three times will scale the eigenvector by $\lambda^3 = (-3)^3 = -27$.
    Hence,
    \begin{align}
        \Vec{M}^3\myvec{1\\-\tfrac{1}{2}\\0}
        = (-27)\myvec{1\\-\tfrac{1}{2}\\0}
        = \myvec{-27\\\tfrac{27}{2}\\0}
    \end{align}
    Thus,
    \begin{align}
        \boxed{|\alpha| = 27.}
    \end{align}
\end{frame}




\end{document}
